Los resultados obtenidos a lo largo del trabajo permiten reinterpretar el diseño de circuitos digitales desde una perspectiva fuertemente dependiente de las no idealidades del nivel transistor. Lejos de ser un aspecto secundario, el dimensionamiento de los dispositivos CMOS condiciona de manera crítica el punto de operación de las compuertas, introduciendo asimetrías que impactan directamente en la robustez lógica y en la tolerancia al ruido.

Asimismo, el análisis dinámico pone de manifiesto que los efectos capacitivos y la interacción entre etapas no pueden ser despreciados, ya que determinan el desempeño temporal del sistema y pueden amplificar desbalances presentes en etapas previas. Esto evidencia que el diseño en cascada no es simplemente modular, sino que requiere una consideración global del circuito.

En el caso de las estructuras secuenciales, se observa que las decisiones arquitectónicas, como el uso de \textit{transmission gates}, no solo optimizan el funcionamiento, sino que resultan determinantes para garantizar la integridad de la información almacenada frente a condiciones no ideales.

En definitiva, el trabajo demuestra que el comportamiento macroscópico de los circuitos digitales está intrínsecamente ligado a fenómenos de bajo nivel, y que un diseño preciso a escala de transistor es indispensable para lograr sistemas confiables y predecibles.